% Artículo VII. Incorporación de resultados preexistentes, 10-09-2026.
% Propietarios íntegros preservados en fuentes_conservadas:
% 17_pantalla_cubica_soldadura.tex (P), 18_clausura_energetica.tex (E).
% P: incidencia, cociente, soldadura celular y refinamiento.
% E: respuesta, unicidad, extensión mínima y expectativa tracial.
% Dependencias anteriores de VII: núcleo APP--TRIT--TPK; realización
% cúbica de la ruta, cociclo de transporte y compresión q_vis=5/9.
% Este fragmento no constituye por sí solo el artículo autónomo.
% La identidad de pantalla se distingue de la contracción material de espín.

\section{Soldadura de frontera y respuesta energética}
\label{vii:sec:pantalla-respuesta}

La realización cúbica del transporte conserva seis incidencias orientadas.
Su cociente geométrico, la forma lorentziana y la respuesta energética
se obtienen de la misma matriz de incidencias. El transporte procede de
la ruta enriquecida APP--TRIT--TPK: la cara, el signo, el acarreo y la
memoria acompañan a cada arista. En esta sección se desarrolla la
operación lineal que actúa sobre esa representación de frontera.

\subsection{Incidencia cúbica y cociente de rango cuatro}
\label{vii:sec:cociente-pantalla}

Sea \(H_F=\mathbb R^F\), con
\(F=\{(i,\epsilon):1\leq i\leq3,\ \epsilon\in\{-1,1\}\}\),
su producto escalar euclidiano y su base \(e_{i,\epsilon}\).
La oposición de caras es la involución ortogonal
\(Re_{i,\epsilon}=e_{i,-\epsilon}\). Definimos
\[
 u=\sum_{i,\epsilon}e_{i,\epsilon},\qquad
 P_u=\frac{uu^{\mathsf T}}6,\qquad P_-=\frac{I-R}{2},\qquad
 P_0=\frac{I+R}{2}-P_u,\qquad P_\square=P_u+P_-.
\]
Los vértices son \(V=\{-1,1\}^3\). La matriz de incidencia
\(M:\mathbb R^V\longrightarrow H_F\) queda determinada por
\[
 M_{(i,\epsilon),\nu}=\mathbf1_{\{\nu_i=\epsilon\}}.
\]

\begin{proposition}[Descomposición de la incidencia]
\label{vii:prop:incidencia}
Se tiene
\[
 MM^{\mathsf T}=12P_u+4P_-,
 \qquad
 \ker M^{\mathsf T}=P_0H_F.
\]
Por consiguiente, el cociente
\(Q=H_F/P_0H_F\), identificado ortogonalmente con
\(P_\square H_F\), tiene dimensión cuatro.
\end{proposition}

\begin{proof}
Una cara contiene cuatro vértices. Dos caras opuestas carecen de vértices
comunes y dos caras de ejes diferentes tienen dos vértices comunes.
Escribiendo \(\mathbf J=uu^{\mathsf T}\), resulta
\[
 MM^{\mathsf T}=4I+2(\mathbf J-I-R)
 =2(I-R)+2\mathbf J=4P_-+12P_u.
\]
Los tres proyectores \(P_u,P_-,P_0\) son ortogonales y suman \(I\);
sus rangos son, respectivamente, uno, tres y dos. La identidad
\(\|M^{\mathsf T}x\|^2=\langle x,MM^{\mathsf T}x\rangle\)
identifica el núcleo anunciado. El rango complementario es cuatro.
\end{proof}

Sobre \(Q\), la forma
\[
 B=P_--P_u
\]
tiene signatura \((3,1)\). La orientación temporal se fija por \(u\).
Con
\[
 t=\frac{u}{\sqrt6},\qquad
 d_i=\frac{e_{i,+}-e_{i,-}}{\sqrt2},\qquad
 W=\sqrt6P_u+\sqrt2P_-,
\]
se cumple \(We_{i,\epsilon}=t+\epsilon d_i\). Esos seis vectores
son nulos para \(B\), mientras que
\[
 3t+\nu_1d_1+\nu_2d_2+\nu_3d_3
\]
tiene norma cuadrática \(-6\). La misma incidencia satisface
\(MM^{\mathsf T}=2W^*W\) sobre \(Q\). Las rotaciones propias del cubo
conmutan con los proyectores, preservan ambas formas y fijan \(t\).
Estas identidades determinan la métrica de la representación cúbica.

\subsection{Soldadura celular y transporte de memoria}
\label{vii:sec:soldadura-celular}

Para cada arista orientada \(a\), la representación de la ruta proporciona
la cara \(\lambda(\underline a)\), el signo
\(\sigma(a)\), las fibras de origen y destino y el transporte
\(U(a):Q_{s(a)}\longrightarrow Q_{t(a)}\).
La inversión verifica \(U(\bar a)=U(a)^{-1}\).
Las semiaristas de frontera conservan su inversión formal y su incidencia
en el borde. Con \(\ell_m=\ell_0\,9^{-m}>0\), la soldadura es
\begin{equation}
 \beta_m(a)=\ell_m\sigma_m(a)
              n_{\lambda(\underline a)}^{s(a)},\qquad
 n_{i,\epsilon}=t+\epsilon d_i.
 \label{vii:eq:soldadura-arista}
\end{equation}
La cara y el signo se transportan junto con la memoria de la ruta.
La convención de inversión exige
\begin{equation}
 \beta_m(\bar a)=-U_m(a)\beta_m(a).
 \label{vii:eq:inversion-soldadura}
\end{equation}
Aplicando dos veces esta relación se recupera \(\beta_m(a)\);
el registro de acarreo conserva, por separado, la composición ordenada.

\begin{proposition}[Naturalidad de la soldadura por refinamiento]
\label{vii:prop:refinamiento-soldadura}
Sea \(a=a_1\cdots a_9\) un refinamiento compatible de una arista.
Supóngase que las nueve incidencias hijas, transportadas a la fibra del
predecesor, conservan la cara y el signo, y que
\(\ell_{m+1}=\ell_m/9\). Sea
\(\mathcal U_j:Q_{m,s(a)}\longrightarrow Q_{m+1,s(a_j)}\)
el transporte acumulado que incluye la identificación por refinamiento
y el trayecto hasta el origen del sucesor \(j\). Entonces
\[
 \beta_m(a)=
 \sum_{j=1}^{9}\mathcal U_j^{-1}\beta_{m+1}(a_j).
\]
La igualdad es compatible con la inversión y con la concatenación
del registro de memoria.
\end{proposition}

\begin{proof}
Cada sumando, expresado en la fibra del origen del predecesor, es
\((\ell_m/9)\sigma_m(a)n_{\lambda(\underline a)}\).
La suma de nueve términos reproduce la soldadura del predecesor.
La inversión invierte el orden de los transportes y sustituye cada uno
por su inverso; la relación \eqref{vii:eq:inversion-soldadura} proporciona
el signo global. La memoria se concatena en ese mismo orden, por lo
que el refinamiento conserva tanto la representación vectorial como su
registro enriquecido.
\end{proof}

\subsection{Respuesta positiva, unicidad y extensión mínima}
\label{vii:sec:respuesta-positiva}

La soldadura lineal de pantalla es
\(\bar\beta_m=\ell_mW|_Q\). Es invertible, con
\[
 \bar\beta_m^{-1}
 =\ell_m^{-1}\left(\frac1{\sqrt6}P_u+\frac1{\sqrt2}P_-\right).
\]
El autovalor \(5/9\) del sector antisimétrico del bloque normalizado
\(B_c/9\), construido en la sección~\ref{vii:sec:bloque-compresion},
determina el coeficiente de
compresión \(q_{\mathrm{vis}}=5/9\). Reservamos esta notación para
distinguirlo del cociente euclídeo \(q_9(n)\) del núcleo aritmético.
Su realización sobre la pantalla utiliza una isometría
\(V_9:Q\longrightarrow\mathcal K\). La contribución visible
\(C_9=q_{\mathrm{vis}}V_9\) y su complemento
\(D_9=\sqrt{1-q_{\mathrm{vis}}^2}\,V_9\) satisfacen
\[
 C_9^*C_9+D_9^*D_9=I_Q,\qquad D_9^*D_9=\frac{56}{81}I_Q.
\]
El coeficiente \(56\) es \(9^2-5^2\); pertenece al defecto cuadrático
de ese lector de compresión.

\begin{theorem}[Respuesta energética determinada por la soldadura]
\label{vii:thm:respuesta-unica}
Existe un único operador \(\Lambda_m:Q\longrightarrow Q\) tal que
\(\bar\beta_m^*\Lambda_m\bar\beta_m=D_9^*D_9\). Es positivo definido y
\begin{equation}
 \Lambda_m=\frac{28}{243\ell_m^2}(P_u+3P_-)
 =\frac{112}{81\ell_m^2}
       (MM^{\mathsf T}|_Q)^{-1}.
 \label{vii:eq:lambda-pantalla}
\end{equation}
Sobre \(H_F\), la identidad correspondiente tiene como segundo miembro
\((56/81)P_\square\).
\end{theorem}

\begin{proof}
La congruencia por la soldadura invertible determina necesariamente
\[
 \Lambda_m=\bar\beta_m^{-*}\frac{56}{81}I_Q\bar\beta_m^{-1}.
\]
Sustituyendo la inversa, el coeficiente de \(P_u\) es
\(56/(486\ell_m^2)=28/(243\ell_m^2)\); el de \(P_-\) es
\(28/(81\ell_m^2)\). Ambos son positivos. Por la
proposición \ref{vii:prop:incidencia},
\((MM^{\mathsf T}|_Q)^{-1}=P_u/12+P_-/4\),
lo que da la segunda expresión. La extensión a seis caras anula
\(P_0H_F\); por ello la identidad se extiende mediante \(P_\square\).
\end{proof}

\begin{proposition}[Caracterización por mínima extensión]
\label{vii:prop:minima-extension}
Para todo \(b\in Q\),
\[
 \langle b,\Lambda_m b\rangle
 =\frac{112}{81\ell_m^2}
   \min_{Mx=b}\|x\|^2.
\]
La extensión minimizante es única:
\(x_{\min}=M^{\mathsf T}(MM^{\mathsf T}|_Q)^{-1}b\).
\end{proposition}

\begin{proof}
El vector anunciado satisface \(Mx_{\min}=b\) y pertenece a
\(\operatorname{im}M^{\mathsf T}=(\ker M)^\perp\).
Toda otra solución es \(x_{\min}+z\), con \(z\in\ker M\); por
ortogonalidad, su norma al cuadrado es \(\|x_{\min}\|^2+\|z\|^2\).
Además,
\[
 \|x_{\min}\|^2
 =\langle b,(MM^{\mathsf T}|_Q)^{-1}b\rangle.
\]
La fórmula de \(\Lambda_m\) concluye la prueba.
\end{proof}

El operador \(MM^{\mathsf T}|_Q\) realiza la respuesta
Dirichlet--a--Neumann de esta formulación de incidencia; su inversa
es la respuesta Neumann--a--Dirichlet. La energía anterior corresponde
a esta última, con la normalización de la soldadura y del defecto de
compresión. El complemento de Schur de
\(\left(\begin{smallmatrix}I&M^{\mathsf T}\\M&0\end{smallmatrix}\right)\)
es \(-MM^{\mathsf T}\), lo que ofrece la misma eliminación variacional.

\subsection{Conservación de la densidad por refinamiento}
\label{vii:sec:densidad-refinamiento}

Consideremos el sector de transportes cúbicos, ortogonales para el producto
escalar energético y compatibles con la forma \(B\). Fijamos un
espacio de etiquetas común \(Q^{\mathrm{lab}}\), identificado
isométricamente con la pantalla del predecesor. Las soldaduras tienen
los tipos
\(\bar\beta_m:Q^{\mathrm{lab}}\longrightarrow Q_m\) y
\(\bar\beta_{m+1,j}:Q^{\mathrm{lab}}\longrightarrow Q_{m+1,j}\);
los operadores \(\Lambda\) actúan en sus respectivas fibras de llegada.
Si
\(U_j:Q_j\longrightarrow Q\) transporta la fibra hija a la del predecesor,
entonces
\[
 \bar\beta_{m+1,j}=\frac19U_j^{-1}\bar\beta_m,\qquad
 \Lambda_{m+1,j}=81U_j^{-1}\Lambda_mU_j.
\]
El factor \(81\) procede de \(\ell_{m+1}^{-2}=81\ell_m^{-2}\).

\begin{theorem}[Invariancia de la densidad normalizada]
\label{vii:thm:densidad}
Cada sucesor conserva la respuesta cuadrática del predecesor:
\[
 \bar\beta_{m+1,j}^*\Lambda_{m+1,j}\bar\beta_{m+1,j}
 =\frac{56}{81}I_{Q^{\mathrm{lab}}}.
\]
Para \(r\) refinamientos, la respuesta conjunta, expresada en las fibras
del predecesor, es
\[
 E_{m+r}=\frac{56}{81}I_{Q^{\mathrm{lab}}}\otimes I_{9^r}.
\]
La expectativa con traza normalizada
\(\tau_{9^r}(A)=9^{-r}\operatorname{Tr}(A)\) satisface
\[
 (\operatorname{id}\otimes\tau_{9^r})(E_{m+r})
 =\frac{56}{81}I_{Q^{\mathrm{lab}}}
\]
a toda profundidad finita.
\end{theorem}

\begin{proof}
La congruencia de un sucesor contiene los factores \(1/9,81,1/9\);
su producto es uno. La ortogonalidad cancela los transportes
intermedios. La suma directa de nueve respuestas iguales identifica
el primer refinamiento con \((56/81)I_Q\otimes I_9\).
La iteración produce el tensor con \(I_{9^r}\).
Finalmente \(\tau_{9^r}(I_{9^r})=1\). Las expectativas se componen
por la identidad
\(\tau_{9^{r+s}}=\tau_{9^r}\otimes\tau_{9^s}\).
\end{proof}

Las inclusiones unitales \(A\mapsto A\otimes I_9\) son compatibles con
esa colección de trazas. En el límite inductivo de las álgebras matriciales,
la respuesta es el elemento central
\((56/81)I_Q\otimes\mathbf1\). La suma extensiva sin normalizar produce,
en cambio, \(9^r(56/81)I_Q\). Esta distinción especifica qué magnitud
se conserva: la densidad cuadrática de pantalla.

\subsection{Realización variacional de la respuesta}
\label{vii:sec:realizacion-respuesta}

La realización material usa una isometría
\(J_{\mathrm{scr}}:Q\longrightarrow\mathcal H_\partial\)
compatible con la acción de espín y con el refinamiento.
La cotetrada de frontera es \(e=J_{\mathrm{scr}}\bar\beta_m\).
Sea \(\Lambda_{\mathrm{CE}}\) la respuesta de frontera obtenida
por segunda variación de la acción material y geométrica, con sus
condiciones de contorno y la eliminación de las variables interiores.
El defecto de compatibilidad es
\begin{equation}
 \mathcal R_{\mathrm{grav}}
 =\bar\beta_m^*
   (J_{\mathrm{scr}}^*\Lambda_{\mathrm{CE}}J_{\mathrm{scr}}-\Lambda_m)
   \bar\beta_m
 =e^*\Lambda_{\mathrm{CE}}e-D_9^*D_9.
 \label{vii:eq:defecto-respuesta}
\end{equation}
Por la invertibilidad de \(\bar\beta_m\),
\[
 \mathcal R_{\mathrm{grav}}=0
 \quad\Longleftrightarrow\quad
 J_{\mathrm{scr}}^*\Lambda_{\mathrm{CE}}J_{\mathrm{scr}}=\Lambda_m.
\]
La isometría \(J_{\mathrm{scr}}\) transporta los estados de frontera;
la corriente variacional de espín es una forma de grado tres con valores
en bivectores. Cada una interviene con su dominio propio en la
realización de la respuesta.

La construcción determina una forma energética positiva, su extensión
interior minimizante y la densidad conservada por refinamiento.
La amplitud de una contribución material concreta se obtiene al evaluar
la corriente de esa realización en su forma cuadrática variacional.
Así quedan diferenciados el operador que fija la respuesta, la corriente
sobre la que actúa y el observable gravitatorio que expresa la evaluación.
